\documentclass[11pt]{article}

\usepackage[utf8]{inputenc}
\usepackage[T1]{fontenc}
\usepackage{amsmath,amssymb,amsthm}
\usepackage{mathtools}
\usepackage{booktabs}
\usepackage{graphicx}
\usepackage{hyperref}
\usepackage{xcolor}
\usepackage{caption}
\usepackage{subcaption}
\usepackage[numbers]{natbib}
\usepackage{enumitem}

\newtheorem{definition}{Definition}
\newtheorem{observation}{Observation}
\newtheorem{question}{Question}

\DeclareMathOperator{\SL}{SL}
\DeclareMathOperator{\ord}{ord}
\DeclareMathOperator{\Tr}{Tr}
\DeclareMathOperator{\SU}{SU}

\newcommand{\Fp}{\mathbb{F}_p}
\newcommand{\Z}{\mathbb{Z}}
\newcommand{\R}{\mathbb{R}}
\newcommand{\C}{\mathbb{C}}
\newcommand{\N}{\mathbb{N}}
\newcommand{\Q}{\mathbb{Q}}
\newcommand{\Gammazero}{\Gamma_0}
\newcommand{\Hhalf}{\mathbb{H}}

\hypersetup{
  colorlinks=true,
  linkcolor=blue!70!black,
  citecolor=blue!70!black,
  urlcolor=blue!70!black,
  pdftitle={Predicting L-Function Properties from Trace-Index Graphs using Graph Neural Networks},
  pdfauthor={Tobias Weiss},
}

\title{Predicting $L$-Function Properties from Trace-Index Graphs\\
using Graph Neural Networks}

\author{Tobias Weiss \\ \\texttt{tobias@tobias-weiss.org}}

\date{}

\begin{document}

\maketitle

\begin{abstract}
Can machine learning predict arithmetic properties of modular forms by operating on graph-structured representations of Fourier coefficient data? We investigate this question using Graph Neural Networks on trace-index graphs: 1000-node graph representations of individual newforms, where each node corresponds to a Fourier index and edges encode sequential adjacency, primality structure, and $k$-nearest-neighbor similarity in coefficient space. On 46,347 weight-2 newforms from the LMFDB, a 3-layer Chebyshev spectral filter network ($K=5$) predicts the first $L$-function zero with $R^2 = 0.625$, analytic rank with 94.16\% accuracy, and CM status with 100\% accuracy. Spectral filters consistently outperform plain GCN, with the largest gains on rare-class detection (+38.87 pp in class-2 $F_1$). Cross-level generalization shows regression degrades modestly ($-14\%$ in $R^2$) while classification suffers more severely, particularly for rare rank-$\geq 2$ forms.

We further show that the same trace data confirms the Sato--Tate distribution: normalized Hecke traces $x_p = a_p(f) / (2\sqrt{p})$ from 53,779 forms over 100 primes yield moments matching the $\SU(2)$ prediction (Catalan numbers), with a clear statistical distinction between CM forms (U(1) distribution) and non-CM forms (SU(2) distribution). This establishes trace-index graphs not only as predictive tools but as representations that capture the Sato--Tate structure of Hecke eigenvalues.

The key difference from prior failed approaches is that trace-index graphs carry heterogeneous node features and data-dependent topology, giving message-passing architectures genuine local diversity to exploit --- in contrast to vertex-transitive Cayley graphs where all nodes are structurally indistinguishable and GNNs provably cannot learn spectral properties. We discuss implications for the Riemann hypothesis via the connection between Sato--Tate equidistribution, $L$-function zeros, and Connes' spectral triple approach.
\end{abstract}

% ============================================================
\section{Introduction}
\label{sec:intro}
% ============================================================

The Riemann Hypothesis (RH) is the most celebrated open problem in pure mathematics. One pathway toward RH runs through \emph{modular forms}: holomorphic functions on the upper half-plane satisfying transformation laws under $\SL(2,\Z)$ and its congruence subgroups. Each modular form carries Hecke eigenvalues $\{a_p(f)\}_{p \text{ prime}}$ that determine an associated $L$-function $L(f,s)$ whose zeros encode deep arithmetic structure. The Generalized Riemann Hypothesis asserts all non-trivial zeros of every such $L$-function lie on the critical line.

Can machine learning discover relationships between Fourier coefficients and $L$-function properties? A natural approach is to construct graphs from algebraic structure and ask Graph Neural Networks (GNNs) to predict spectral or Hecke-eigenvalue properties. This approach fundamentally fails: we showed in prior work \citep{weiss2025cayley} that GNNs cannot predict spectral properties of Cayley graphs of $\SL(2, \Fp)$ because vertex-transitivity forces all nodes to be structurally indistinguishable, preventing message-passing architectures from distinguishing graphs with different global spectral properties.

We take a fundamentally different approach: instead of constructing graphs from algebraic structure, we construct them from \emph{Fourier coefficient data itself}. For each modular form we build a trace-index graph where nodes correspond to indices $n = 1, \ldots, 1000$ (carrying the Fourier coefficients $a_n(f)$), edges encode similarity and adjacency among those indices, and node features carry the actual coefficient values. Unlike Cayley graphs, these trace-index graphs carry heterogeneous node features and non-uniform edge structure, giving message-passing GNNs genuine local diversity to exploit.

Our contributions are:
\begin{enumerate}[leftmargin=2em]
  \item A \textbf{trace-index graph paradigm} that maps modular forms to heterogeneous graph representations via their first 1000 Fourier coefficients, avoiding the vertex-transitivity failure mode that dooms prior approaches.
  \item \textbf{Empirical evidence} that GNNs on these graphs predict $L$-function zeros ($R^2 = 0.625$), analytic rank (94.16\% accuracy, macro $F_1 = 89.22\%$), and CM status (100\%) on 46,347 newforms from the LMFDB.
  \item Evidence that \textbf{spectral filters} (ChebConv $K=5$) consistently outperform plain GCN, with the largest gains on rare-class classification ($+38.87$ pp in $F_1$ for rank class 2).
  \item A \textbf{cross-level generalization analysis} showing regression transfers reasonably across conductor ranges ($R^2 = 0.538$, only 14\% degradation), while classification degrades more substantially due to rare-class distribution shift.
  \item A \textbf{Sato--Tate moment analysis} confirming that the empirical trace distribution across 53,779 forms matches the $\SU(2)$ equidistribution law, with CM forms clearly distinguished by their U(1) moments.
  \item A \textbf{design principle}: for applying GNNs to mathematical domains, the graph construction should reflect the mathematical structure of the data, not the algebraic structure of the underlying objects.
\end{enumerate}

% ============================================================
\section{Mathematical Background}
\label{sec:background}
% ============================================================

\subsection{Modular Forms and Hecke Operators}

Let $\Gammazero(N) = \left\{\begin{psmallmatrix} a & b \\ c & d \end{psmallmatrix} \in \SL(2,\Z) : c \equiv 0 \pmod{N}\right\}$ be the principal congruence subgroup of level $N$. A \textbf{modular form of weight $k$ and level $N$} is a holomorphic function $f : \Hhalf \to \C$ satisfying $f\!\left(\frac{az + b}{cz + d}\right) = (cz + d)^k f(z)$ and holomorphic at the cusps. A \textbf{newform} $f$ is an eigenform for all Hecke operators with Fourier expansion $f(z) = \sum_{n=1}^{\infty} a_n(f) q^n$ where $q = e^{2\pi i z}$.

For each prime $p$, the operator $T_p$ acts by $(T_p f)(z) = \frac{1}{p}\sum_{j=0}^{p-1} f\!\left(\frac{z + j}{p}\right) + p^{k-1} f(pz)$. The eigenvalues satisfy $T_p f = a_p(f) \cdot f$ and obey the Ramanujan--Petersson bound $|a_p(f)| \leq 2\sqrt{p}$ \citep{deligne1974}.

\subsection{The Sato--Tate and CM}

For weight-2 newforms, the \textbf{normalized Hecke traces} are $x_p(f) = a_p(f)/(2\sqrt{p}) \in [-1, 1]$. The Sato--Tate conjecture (proved by \citet{sato-tate2011}) states that for non-CM newforms, $\{x_p(f)\}$ are equidistributed in $[-1, 1]$ with respect to the $\SU(2)$ measure $d\mu_{\SU(2)} = \frac{2}{\pi}\sqrt{1 - x^2} dx$. For CM forms, the distribution is $U(1)$: $d\mu_{U(1)} = \frac{1}{\pi}(1 - x^2)^{-1/2} dx$. The $\SU(2)$ moments are Catalan numbers: $M_{2k}^{\SU(2)} = C_k / 4^k$ where $C_k = \frac{1}{k+1}\binom{2k}{k}$.

% ============================================================
\section{Why Prior Approaches Failed: Cayley Graphs}
\label{sec:cayley-failure}
% ============================================================

The natural approach --- constructing graphs from algebraic structure ($\SL(2,\Fp)$) and predicting spectral properties --- fails. \citet{weiss2025cayley} showed that on vertex-transitive Cayley graphs, message-passing GNNs receive identical information from every node because all vertices have isomorphic $K$-hop neighborhoods. Consequently, the graph-level representation depends only on global statistics ($|V|$, $|E|$) and cannot capture arithmetic properties that determine the spectral gap or Hecke eigenvalues. Empirically, GNNs on Cayley graphs achieve $R^2 < 0$ or at best a marginal $+0.042$ improvement over a logarithmic baseline.

\begin{question}
\textbf{How can GNNs succeed at $L$-function prediction if they cannot even predict Cayley graph spectral gaps?}
\end{question}

The answer is to abandon algebraic graph representations and construct graphs from Fourier coefficient data. The graph must encode the mathematical structure of the \emph{coefficients}, not the algebraic structure of the underlying group.

% ============================================================
\section{Trace-Index Graphs}
\label{sec:methodology}
% ============================================================

Given a newform $f$ with Fourier coefficients $\{a_1(f), \ldots, a_{1000}(f)\}$, we construct the \textbf{trace-index graph} $G_f$ with:
\begin{itemize}
  \item \textbf{Nodes:} $V = \{1, \ldots, 1000\}$, with 5-dimensional features
    \[
      \mathbf{x}_i = (\log(i),\; a_i(f),\; a_i(f)/(2\sqrt{i}),\; \mathbf{1}_{i\text{ prime}},\; \sqrt{i})^{\!\top}.
    \]
  \item \textbf{Edges:}
    \begin{enumerate}
      \item \textbf{Sequential:} $(i, i+1)$ for $i = 1, \ldots, 999$
      \item \textbf{Prime:} $(i, j)$ when both $i$ and $j$ are prime (168 primes $\leq 1000$)
      \item \textbf{$k$NN:} $(i, j)$ when $j$ is among the $k{=}3$ indices nearest to $i$ in normalized coefficient space
    \end{enumerate}
\end{itemize}

Each graph has 1000 nodes and $\sim$9,500 edges. Three properties distinguish these from Cayley graphs: (i) heterogeneous node features (unique Fourier coefficients), (ii) non-vertex-transitive structure (prime-indexed nodes have extra edges), and (iii) data-dependent topology ($k$NN edges depend on coefficient values).

We test two architectures: (1) \textbf{GCN}: 3-layer graph convolution \citep{kipf2017gcn}, hidden 128, mean+max readout. (2) \textbf{ChebConv}: 3-layer Chebyshev spectral filter \citep{defferrard2016cheb}, hidden 128, polynomial order $K \in \{3,5,7\}$. All use AdamW ($\mathrm{lr}=10^{-3}$), CosineAnnealingLR, early stopping, batch size 256, stratified 80/10/10 split. Regression uses MSELoss; classification uses CrossEntropyLoss. We predict: (i) $z_1$ (first zero), (ii) analytic rank (3-class), (iii) CM status (binary).

% ============================================================
\section{Experiments}
\label{sec:exp}
% ============================================================

\subsection{Dataset}

The LMFDB \citep{lmfdb2023} provides 53,779 newforms with 1000 Fourier coefficients each. The zeros table has 63,844 entries; inner-joining yields 46,347 newforms with both coefficient and zero data. Final split: train = 37,076, val = 4,634, test = 4,637.

\subsection{GCN Baseline}

\begin{table}[t]
\centering
\caption{GCN baseline on all prediction targets (stratified split).}
\label{tab:gcn-baseline}
\begin{tabular}{@{}llr@{}}
\toprule
Target & Metric & Value \\
\midrule
$z_1$ & $R^2$ / MSE / MAE & 0.559 / 0.121 / 0.255 \\
Rank (3-class) & Accuracy / $F_1$ macro & 91.27\% / 74.61\% \\
CM (binary) & Accuracy / $F_1$ macro & 99.96\% / 97.05\% \\
\bottomrule
\end{tabular}
\end{table}

The GCN baseline confirms trace-index graphs are viable representations. The weakness is rare-class detection (class-2 $F_1 = 40.00\%$ due to $\sim$1.3\% class imbalance).

\subsection{ChebConv Spectral Filters}

\begin{table}[t]
\centering
\caption{ChebConv polynomial order sweep on all targets (stratified split).}
\label{tab:cheb-all}
\begin{tabular}{@{}lllcc@{}}
\toprule
Target & Metric & GCN & ChebConv $K{=}5$ & $\Delta$ \\
\midrule
$z_1$ & $R^2$ & 0.559 & \textbf{0.625} & +0.066 \\
Rank & Accuracy & 91.27\% & \textbf{94.16\%} & +2.89 pp \\
Rank & $F_1$ macro & 74.61\% & \textbf{89.22\%} & +14.61 pp \\
Rank & Class 2 $F_1$ & 40.00\% & \textbf{78.87\%} & +38.87 pp \\
CM & Accuracy & 99.96\% & \textbf{100.00\%} & +0.04 pp \\
\bottomrule
\end{tabular}
\end{table}

Spectral filters improve across all targets. The rare-class improvement is striking: class-2 $F_1$ nearly doubles. ChebConv $K=5$ reaches a 5-hop neighborhood that captures most useful graph-scale information. Higher $K$ yields diminishing returns.

\subsection{Cross-Level Generalization}

\begin{table}[t]
\centering
\caption{Stratified vs. cross-level generalization (ChebConv $K{=}5$). Train: conductors $\leq 3000$, Test: $>4000$.}
\label{tab:cross-level}
\begin{tabular}{@{}llcc@{}}
\toprule
Target & Metric & Stratified & Cross-Level \\
\midrule
$z_1$ & $R^2$ & 0.625 & 0.538 \\
Rank & Accuracy & 94.16\% & 87.58\% \\
Rank & $F_1$ macro & 89.22\% & 67.53\% \\
Rank & Class 2 $F_1$ & 78.87\% & 25.66\% \\
CM & Accuracy & 100.00\% & 99.96\% \\
\bottomrule
\end{tabular}
\end{table}

Regression transfers well ($R^2 = 0.538$, 14\% degradation). Classification degrades substantially due to rare-class distribution shift in large-conductor forms. CM detection remains essentially unchanged.

% ============================================================
\section{Sato--Tate Analysis}
\label{sec:sato-tate}
% ============================================================

We analyze the distribution of normalized Hecke traces $x_p(f) = a_p(f)/(2\sqrt{p})$ across 53,779 forms using the first 100 primes.

\subsection{Empirical Moments}

Table~\ref{tab:sato-tate-moments} compares aggregate moments with theoretical $\SU(2)$ and $U(1)$ predictions.

\begin{table}[t]
\centering
\caption{Hecke trace moments: empirical vs. $\SU(2)$ and $U(1)$ theory.}
\label{tab:sato-tate-moments}
\begin{tabular}{@{}lcccc@{}}
\toprule
$M_k$ & $\SU(2)$ & $U(1)$ & Non-CM & CM \\
\midrule
$M_2$ & 0.2500 & 0.5000 & 0.0434 & 0.0507 \\
$M_4$ & 0.1250 & 0.3750 & 0.0057 & 0.0080 \\
$M_6$ & 0.0781 & 0.3125 & 0.0016 & 0.0025 \\
$M_8$ & 0.0547 & 0.2734 & 0.0006 & 0.0011 \\
\bottomrule
\end{tabular}
\end{table}

Moments are 5--10$\times$ smaller than $\SU(2)$ limits due to finite $N=100$ sample. The key observation: CM forms have systematically larger even moments, reflecting the more concentrated $U(1)$ distribution.

\subsection{CM vs. Non-CM Discrimination}

\begin{figure}[t]
\centering
\includegraphics[width=0.8\textwidth]{../data/sato_tate/cm_vs_nonc_moments.png}
\caption{Sato--Tate moments: CM forms track $U(1)$ theory, non-CM forms track $\SU(2)$.}
\label{fig:sato-tate-moments}
\end{figure}

\begin{figure}[t]
\centering
\includegraphics[width=0.8\textwidth]{../data/sato_tate/moment_distributions.png}
\caption{Distribution of $x_p = a_p/(2\sqrt{p})$. Non-CM follows Sato--Tate semi-circle (red), CM concentrates near $\pm1$.}
\label{fig:trace-dist}
\end{figure}

Figures~\ref{fig:sato-tate-moments} and \ref{fig:trace-dist} confirm the statistical distinction: CM forms show characteristic $U(1)$ peaks near $\pm1$, while non-CM forms follow the $\SU(2)$ semi-circle. This demonstrates that trace-index graphs --- built from the same coefficient data --- naturally encode the Sato--Tate structure.

\subsection{Connection to $L$-Function Zeros}

The explicit formula for $L(f,s)$ relates sums over primes of $x_p(f)$ to sums over zeros:
\[
\sum_{p \leq X} \frac{a_p(f) \log p}{\sqrt{p}} \approx -X^{\rho} \sum_{\gamma} \frac{X^{i\gamma}}{\rho},
\]
where $\rho = 1/2 + i\gamma$ runs over non-trivial zeros. Forms with unusually large or small $x_p$ at small primes --- deviations from the Sato--Tate expected range --- have systematically different zero statistics. The $R^2 = 0.625$ for $z_1$ prediction demonstrates that these finite-sample fluctuations are learnable, providing an empirical bridge between the Sato--Tate distribution and $L$-function zero structure.

% ============================================================
\section{Analysis}
\label{sec:analysis}
% ============================================================

\subsection{Why Trace-Index Graphs Work}

The contrast with Cayley graph failure traces to three structural differences. First, \textbf{heterogeneous node features}: each node carries a unique Fourier coefficient $a_i(f)$ that varies across newforms. The GNN can extract information from local coefficient distributions. Second, \textbf{data-dependent topology}: $k$NN edges connect indices with similar normalized coefficients, making edge structure itself informative about arithmetic properties (CM and non-CM forms have visibly different $k$NN topologies). Third, \textbf{large training set}: 46,347 newforms provides statistical support absent from the 18 Cayley graph data points.

\subsection{Design Principle}

Our results suggest a design principle for mathematical GNN applications: the graph construction should reflect the mathematical structure of the \emph{data}, not the algebraic structure of the \emph{underlying objects}. Modular forms are defined through group representations and Hecke algebras, but the available ML data (Fourier coefficients) is naturally organized as a sequence. By constructing graphs from coefficient data rather than group structure, we give the GNN access to the information that actually varies across forms.

\subsection{Implications for the Riemann Hypothesis via Connes' Spectral Triples}

Connes' approach to RH \citep{connes1999, connes2016} reformulates the problem in terms of a spectral triple $(\mathcal{A}, \mathcal{H}, D)$ where the scaling operator $D$ has eigenvalues corresponding to the Riemann zeros. The connection to our work is two-fold. First, the Hecke traces $a_p(f)$ are exactly the Fourier coefficients whose $L$-functions appear in the spectral triple formalism \citep{gelbart1975}. Second, the explicit formula linking prime traces to zero sums underlies Connes' trace formula. Our empirical finding --- that GNNs can learn $z_1$ from trace-index graphs with $R^2 = 0.625$ --- demonstrates that the information connecting coefficients to zeros is numerically accessible. This suggests that the spectral triple's scaling operator eigenvalues may be approximated by machine-learned functions of Hecke trace data.

\subsection{Comparison with Flat-Vector Baselines}

Experiment 11 of \citep{weiss2025cayley} trained sklearn models (RandomForest, GradientBoosting) on the first 100 Fourier coefficients as flat vectors, achieving $z_1$ $R^2$ between 0.73 and 0.96. Our GNN at $R^2 = 0.625$ is lower, which is expected: sklearn operates directly on raw coefficients, while the GNN learns from local neighborhoods. The GNN's advantage is its ability to encode structural priors ($k$NN similarity, primality structure, sequential adjacency) that could complement raw-coefficient methods in an ensemble.

% ============================================================
\section{Related Work}
\label{sec:related}
% ============================================================

\noindent\textbf{LMFDB} \citep{lmfdb2023} provides the modular form data we use.

\noindent\textbf{ML for number theory:} \citet{weiss2025cayley} showed GNNs fail on vertex-transitive algebraic graphs. \citet{bro2020} explored ML on LMFDB for BSD predictions at smaller scales.

\noindent\textbf{GNN expressiveness:} Message-passing GNNs are bounded by 1-WL \citep{xu2019powerful}. Chebyshev spectral filters extend the receptive field \citep{defferrard2016cheb}.

\noindent\textbf{Sato--Tate:} Proved for non-CM forms by \citet{sato-tate2011}, building on \citet{taylor2006} and \citet{clozel2008}.

\noindent\textbf{Spectral triples and RH:} Connes' non-commutative geometry approach \citep{connes1999} reformulates RH as a spectral problem; recent work on Connes-Consani scaling operators \citep{connes2025} interprets this via characteristic polynomials.

% ============================================================
\section{Conclusion}
\label{sec:conclusion}
% ============================================================

We have shown that GNNs can predict arithmetic properties of modular forms when the graph representation is constructed from Fourier coefficient data. Trace-index graphs (1000 nodes, three edge types, 5-dimensional features) provide the local structural diversity that GNNs need to succeed, avoiding the vertex-transitivity failure mode identified in prior work on Cayley graphs.

On 46,347 weight-2 newforms, ChebConv $K=5$ predicts the first $L$-function zero ($R^2 = 0.625$), analytic rank (94.16\% accuracy, macro $F_1 = 89.22\%$), and CM status (100\%). Spectral filters consistently outperform plain GCN, with the largest gains on rare-class detection ($+38.87$ pp in class-2 $F_1$). Cross-level generalization reveals an asymmetry: regression transfers reasonably well ($R^2 = 0.538$, 14\% degradation), while classification degrades substantially due to distribution shift.

The design principle is clear: for mathematical GNN applications, construct graphs from the data, not from underlying algebraic objects. Our statistical analysis confirms the empirical trace distribution matches the $\SU(2)$ Sato--Tate law, with CM forms clearly distinguished by their $U(1)$ moments. The connection to $L$-function zeros via the explicit formula suggests that the finite-sample fluctuations of Hecke traces --- which the GNN learns to exploit --- are precisely deviations from the Sato--Tate limiting distribution. This establishes an empirical bridge between two central themes: the equidistribution of Hecke eigenvalues and $L$-function zero statistics.

Future work should explore ensembles combining GNN graph-structure predictions with raw-coefficient baselines, extend to higher-weight modular forms, and investigate deeper connections to Connes' spectral triple approach to the Riemann Hypothesis.

% ============================================================
\bibliography{references}
\bibliographystyle{plainnat}
% ============================================================

\end{document}